\documentclass[11pt]{article}
\usepackage{fontspec}
\usepackage[french]{babel}
\usepackage{amsmath}
\usepackage{amssymb}
\usepackage{amsthm}
\usepackage{mathrsfs}
\usepackage{tikz-cd}
\usepackage{enumitem}
\usetikzlibrary{babel,decorations.pathreplacing}

\newcommand{\printedpage}[1]{\marginpar{\footnotesize p.~#1}}
\newcommand{\msp}[1]{\marginpar{\footnotesize [#1]}}
\newcommand{\M}{\mathcal{M}}
\newcommand{\Mc}{\mathcal{M}_{\mathrm{c}}}

\begin{document}

\printedpage{19}
% head: A. Grothendieck : Dérivateurs, Ch. XVII, Édition des Universités de Montpellier II et Paris VII

\[
\begin{tikzcd}
X_{I_\alpha} \arrow[r] \arrow[d] & X_{I_{\alpha_0}} \arrow[d] \\
Y_{I_\alpha} \arrow[r] & Y_{I_\alpha,I_{\alpha_0}} \rlap{${}= Z_\alpha$} \arrow[d] \\
& Y_{I_{\alpha_0}} \rlap{\ ,}
\end{tikzcd}
\]
on a $Z_0 = Y_{I_\alpha,I_{\alpha_0}}$, $Z_{\alpha_0} = Y_{I_\alpha,I_{\alpha_0}}$, $Z_\alpha = \varinjlim\limits_{\beta<\alpha} Z_\beta$ si $\alpha$ \msp{34} ordinal limite. Donc il reste à déterminer les $v_{\alpha,\alpha+1} : Z_\alpha \longrightarrow Z_{\alpha+1}$ (pour $\alpha < \alpha_0$), ce qu'on voit sur le diagramme
\[
\begin{tikzcd}
X_{I_\alpha} \arrow[r] \arrow[d] & X_{I_{\alpha+1}} \arrow[r] \arrow[d] & X_{I_{\alpha_0}} \rlap{${}= Z_0$} \arrow[d] \\
Y_{I_\alpha} \arrow[r] & Y_{I_\alpha,I_{\alpha+1}} \arrow[r] \arrow[d, "u_{I_\alpha,I_{\alpha+1}}"] & Y_{I_\alpha,I_{\alpha_0}} \rlap{${}= Z_\alpha$} \arrow[d, "v_{\alpha,\alpha+1}"] \\
& Y_{I_{\alpha+1}} \arrow[r] & Y_{I_{\alpha+1},I_{\alpha_0}} \rlap{${}= Z_{\alpha+1}\ .$}
\end{tikzcd}
\]

Reprenons la démonstration de 2.1. On va prouver 1° par les deux lemmes précédents. Notons que lorsque l'espace sous-jacent à $I$ est noethérien, alors la construction transfinie est forcément finie, donc on n'a pas besoin de la stabilité de $C$ par composés transfinis infinis.

\msp{35}

Prouvons 2°. On suppose donc que pour tout $i_0 \in I$, on ait $u_{I_{<\alpha_0},I/i_0} \in C$, on veut en déduire que $u_\bullet \in C_I$, i.e. $u_{J,I/i_0} \in C$ pour \textit{tout} ouvert.

\noindent\textbf{Lemme 2.8.} (Sous hypothèse artinienne sur $I$.) \textit{Soit $P \subset I$ l'ensemble des $i_0 \in I$ tels que pour $u_\bullet|I/i_0$ soit dans $C_{I/i_0}$, i.e. que pour tout $i \leq i_0$, et tout ouvert $K$ de $I/i$, on ait $u_{K,I/i} \in C$. Soit $i_0 \in I$ tel que $I_{<i_0} \subset P$. Alors $i_0 \in P$.}

\noindent\textsc{Démonstration}. Il reste à prouver que si $K \subset I/i_0 = J$ est un ouvert, alors $u_{K,J} \in C$. On peut supposer $K \neq I/i_0$, i.e. $K \subset I_{<i_0} \overset{\text{déf}}{=} J'$. Mais alors par hypothèse initiale sur $u_\bullet$, on a $u_{J',J} \in C$, et d'autre part, par hypothèse sur $i_0$, $u_\bullet|J'$ est dans $C_{J'}$. Donc par le théorème 2.1. 1°, on a $u_{K,J'} \in C$. Donc par transitivité $u_{K,J} \in C$, q.e.d.

\noindent\textbf{Corollaire 2.9.} \textit{Si $I$ est artinien, alors $P = I$, i.e. $u_\bullet \in C_I$.}

Cela prouve donc 2.1. 2°. Notons que $\mathrm{Esp}(I)$ est noethérien, il en est de même de tout ouvert (tel $J'$ ci-dessus), donc \msp{36} pour la conclusion du lemme, inutile de supposer $C$ stable par composition infinie.

\printedpage{20}
% head: A. Grothendieck : Dérivateurs, Ch. XVII, Édition des Universités de Montpellier II et Paris VII

Prouvons 2.1. 3°. On a a) $\Longrightarrow$ b), car par 2.1. 1°, $u_\bullet \in (TC)_I$ implique que pour $i \in I$, $u_i : X_i \longrightarrow Y_i$ est dans $TC$ (cf. cor. 2.1 a). Donc il reste à prouver b) $\Longrightarrow$ a), i.e. si $u_\bullet \in C \cap W_{I^\circ}$, alors $u_\bullet \in (TC)_I$, i.e. les $u_{K,I/i_0}$ sont dans $TC$. On sait déjà que $u_{K,I/i_0} \in C$ par hypothèse $u_\bullet \in C_I$, donc il faut prouver que c'est dans $W$. On est ramené à prouver que $u_K \in W$, i.e. $u_K \in TC$, car par le diagramme (où $J = I/i_0$)
\[
\begin{tikzcd}
X_K \arrow[r] \arrow[d, "u_K"'] \arrow[dr, phantom, "\text{cocart.}"] & X_J \arrow[d, "\tilde{u}_{K,J}"'] \arrow[dd, bend left=50, "u_J"] \\
Y_K \arrow[r] & Y_{K,J} \arrow[d, "u_{K,J}"'] \\
& Y_J \rlap{\ ,}
\end{tikzcd}
\]
$u_K \in TC$ implique $\tilde{u}_{K,J} \in TC$, et comme $u_J \simeq u_{i_0} \in W$ par hypothèse, il en résulte bien $u_{K,J} \in W$ par saturation. Mais on va s'y prendre autrement :

\msp{37}

\noindent\textbf{Lemme 2.9.} \textit{$I$ ensemble ordonné quelconque, $C$, $W$ comme dans th.} 2.1, 1° \textit{et} 3°, \textit{$u_\bullet : X_\bullet \longrightarrow Y_\bullet$ dans $C_I \cap W_{I^\circ}$. Soit $P$ l'ensemble des $i_0 \in I$ tels que $u_\bullet|I/i_0$ soit dans $(TC)_{I/i_0}$. Si $i_0 \in I$ est tel que $I_{<i_0} \subset P$, alors $i_0 \in P$.}

En effet, il reste à prouver que pour tout ouvert $K$ de $J = I/i_0$, $K \neq J$, on a $u_{K,J} \in TC$. Mais pour ceci, il suffit de voir que $u_K, u_J \in W$ (comme on venait de le voir). Or $u_J \simeq u_{i_0} \in W$ par hypothèse. D'autre part $u_\bullet|K \in (TC_K)$ puisque $K \subset I_{<i_0} \subset P$, donc par le théorème 2.1, 1°, on conclut que $u_K \in TC$.

Si maintenant on suppose $I$ artinien, il résulte de 3.9 que $P = I$, donc $u_\bullet \in (TC)_I$, q.e.d. % [sic] 3.9 -> 2.9, the lemma just stated is Lemme 2.9

Cela achève la démonstration du théorème 2.1, et aussi du corollaire 2.5, vu que pour la démonstration du lemme 2.9 également, les hypothèses infinies sur $C$, $C \cap W$ sont inutiles si $\mathrm{Top}(I)$ est noethérien, donc ses ouverts sont aussi noethériens.

\noindent\textbf{Remarques.}
\begin{enumerate}[leftmargin=2.5em]
\item[1)] Dans 2.1. 2°, l'hypothèse $I$ artinien n'est par superflue, comme on voit en prenant $I = \mathbf{N}^{\mathrm{opp}} \simeq \mathbf{Z}^{\leq 0}$, \msp{38} et pour une flèche % [sic] n'est par superflue -> n'est pas superflue, negation "pas" intended
\[ \emptyset_{\M(I)} \longrightarrow Y_\bullet\ . \]

Dire que cette flèche est dans $C_I$ signifie que les morphismes de transition $Y_i \xrightarrow{u_i} Y_{i+1}$ sont dans $C$, \textit{et} que les $Y_i$ sont $C$-cofibrants, i.e. $\emptyset_{\M} \longrightarrow Y_i$ est dans $C$, pour tout $i$. Mais la condition envisagée dans 2.1. 1° est seulement que les $u_i$ sont dans $C$. Donc en prenant pour $Y_\bullet$ un système inductif constant de valeur $Y$, avec $Y$ non $C$-cofibrant, i.e. $\emptyset_{\M} \longrightarrow Y$ non dans $C$, on a un contre-exemple.

Avec le même $I$, on voit que $u_\bullet : X_\bullet \longrightarrow Y_\bullet$ est dans $C$ si et seulement si
\begin{enumerate}[leftmargin=2.5em]
\item[a)] les $u_i \in C$, et
\item[b)] les $u_{i,i+1} : Y_{i,i+1} \longrightarrow Y_{i+1}$ sont dans $C$.
\end{enumerate}

\printedpage{21}
% head: A. Grothendieck : Dérivateurs, Ch. XVII, Édition des Universités de Montpellier II et Paris VII

Si de plus les $u_i$ sont dans $W$, donc dans $TC$, on voit comme ci-dessus que $u_{i,i+1} \in TC$ pour tout $i$, donc $u_\bullet \in (TC)_I$. Donc ici on a encore
\[ (TC)_I \ = \ C_I \cap W_{I^\circ}\ , \]
quoique $I$ ne soit pas artinien. J'ignore si l'hypothèse artinienne dans 2.1. 3° est superflue ou non.
\item[2)] Soit $u_\bullet \in \mathrm{Fl}(\M(I^\circ))$\footnote[13]{$I^\circ = I^{\mathrm{opp}}$.}, donc $u_\bullet$ est un objet de $\underline{\mathrm{Fl}}(\M(I^\circ)) \simeq \M(I^\circ \times \boldsymbol{\Delta}^1) = \M((I \times (\boldsymbol{\Delta}^1)^{\mathrm{opp}})^{\mathrm{opp}})$. On vérifie que pour que cet objet soit $C$-cofibrant, il faut et il suffit que
\begin{enumerate}[leftmargin=2.5em]
\item[1°)] $X_\bullet$ soit $C$-cofibrant, et
\item[2°)] $u_\bullet \in C_I$ (ce qui implique que $Y_\bullet$ est aussi cofibrant).
\end{enumerate}
\end{enumerate}
\msp{39}

\section*{3\quad Catégories de modèles à cofibrations}

Soit $\M$ une catégorie munie de deux ensembles de flèches
\[ W,\ C\ \subset\ \mathrm{Fl}(\M)\ . \]
On suppose que $\M$ contient un objet initial, et :
\begin{description}[leftmargin=3em, labelwidth=2.5em]
\item[Cof$_0$] $W$ est modérément saturé, i.e. contient les isomorphismes et si $f$, $g$ sont composables, et si deux parmi $f$, $g$, $gf$ sont dans $W$, le troisième aussi.
\item[Cof$_1$] $C$ contient les isomorphismes, est stable par composition et par cochangement de base.
\item[Cof$_2$] $TC \overset{\text{déf}}{=} C \cap W$ est stable par cochangement de base.
\item[Cof$_3$] Toute flèche $f$ dans $\M$ s'écrit
\[ f \ = \ pi \]
avec $i \in C$, $p \in W$.
\end{description}
Ce sont les axiomes de K. S. \textsc{Brown} (forme duale), qui suppose de plus que les objets de $\M$ sont $C$-cofibrants, i.e. $\emptyset_{\M} \longrightarrow X$ est dans $C$. Je ne ferai pas cette hypothèse. La sous-catégorie $\Mc$ de $\M$ formée des \msp{40} objets $C$-cofibrants satisfait aux axiomes de \textsc{Brown}, de plus on prouve que
\[ \mathrm{Ho}_W(\Mc) \xrightarrow{\ \sim\ } \mathrm{Ho}_W(\M) \]
est une équivalence. Ainsi les résultats de \textsc{Brown} relatifs à $\Mc$, muni des ensembles de flèches induits par $W$, $C$, sont applicables à $\mathrm{Ho}_W(\M)$ et à ses avatars (tels $\mathrm{Ho}_W(S\backslash\M)$, $\mathrm{Ho}_W(\M/T)$ et $\mathrm{Ho}_W(S\backslash\underset{u_0}{\M}/T)$, \ où $u_0 : S \longrightarrow T$ est donné). Il se trouve que pour l'essentiel, les résultats homotopiques standard (notamment suite exacte de cofibrations relatifs aux opérations homotopiques “à droite” (ou faut-il dire à gauche ? il s'agit pourtant de questions d'exactitude à droite, non à gauche)) sont valables. Il me faudra revenir dessus de façon plus circonstanciée par la suite. Notamment, il y a une bonne théorie de carrés $W$-cocartésiens (et \textit{dans} $\Mc$\textit{,} il y a assez de $W$-cofibrations\footnote[14]{Plus précisément, par K. S. \textsc{Brown}, \textit{dans} $\Mc$\textit{,} les objets de $C$ sont des $W_{\mathrm{c}}$-cofibrations.}, \ldots),
\printedpage{22}
% head: A. Grothendieck : Dérivateurs, Ch. XVII, Édition des Universités de Montpellier II et Paris VII
\msp{41}
et de plus, ceux-ci sont \textit{exacts.} La catégorie $\mathrm{Ho}_W(\M) \simeq \mathrm{Ho}_W(\Mc)$ “se calcule” par un calcul de fractions à gauche dans $\pi\Mc$, où $\pi\Mc$ est la catégorie déduite de $\Mc$ en y remplaçant les $\mathrm{Hom}_{\Mc}(X,Y) = \mathrm{Hom}_{\M}(X,Y)$ par le quotient par une relation d'homotopie convenable :
\begin{description}[leftmargin=9em, labelwidth=8em, font=\normalfont]
\item[$f,g : X \rightrightarrows Y$] sont QB-homotopes (au sens de \textit{Quillen-Brown}) s'il existe un $X' \xrightarrow{\alpha} X$ dans $W$ (et on peut supposer $X' \in \mathrm{Ob}\,\Mc$ également) tel que $f\alpha \overset{1}{\sim} g\alpha$ au sens de \textsc{Quillen}, i.e. il existe un objet cylindre
\begin{tikzcd}[baseline=(current bounding box.center)]
I_{X'} \arrow[r, "\in W"] & X' \\
X' \vee X' \arrow[u, "\in C"] \arrow[ur, dotted] &
\end{tikzcd}
pour $X'$, tel que $(f\alpha, g\alpha)\ :\ X' \vee X' \longrightarrow Y$ se factorise par un $h' : I_{X'} \longrightarrow Y$.
\end{description}

C'est de ce critère que résulte l'exactitude (à gauche) des carrés $W$-cocartésiens.

Je veux voir comment une structure de cofibrations sur $\M$ s'étend aux catégories $\M(I^\circ) = \underline{\mathrm{Hom}}(I, M)$. On a dans $\M(I^\circ)$ \msp{42} l'ensemble $W_{I^\circ}$, et $C_I$ (il vaut mieux sans doute le noter aussi $C_{I^\circ}$) défini au paragraphe précédent. Il s'agit de vérifier les axiomes. L'axiome $\mathrm{Cof}_0$ est immédiat, et bien connu pour $I$ une catégorie quelconque. % [sic] \underline{Hom}(I, M) -> \underline{Hom}(I, \mathcal{M}), italic M printed where the script \mathcal{M} is used everywhere else

$\mathrm{Cof}_1$ : $C_{I^\circ}$ contient les isomorphismes, c'est clair, parce que $C$ les contient. Stable par composition $X_\bullet \xrightarrow{u} Y_\bullet \xrightarrow{v} Z_\bullet$, regardons $K \subset J$, posons $w = v \circ u$,
\[
\begin{tikzcd}[execute at end picture={\draw[decorate,decoration={brace,mirror,amplitude=6pt}] ([yshift=-2pt]B.south west) -- ([yshift=-2pt]E.south east) node[midway,below=7pt]{$\scriptstyle w_{K,J}$};}]
X_K \arrow[r] \arrow[d] & X_J \arrow[d] & & \\
Y_K \arrow[r] \arrow[d] & Y_{K,J} \arrow[r, "u_{K,J}"] \arrow[d] & Y_J \arrow[d] & \\
Z_K \arrow[r] & |[alias=B]| \bullet \arrow[r, "u'_{K,J}"] & Z_{K,J} \arrow[r, "v_{K,J}"] & |[alias=E]| Z_J \rlap{\ ,}
\end{tikzcd}
\]
\vspace{1ex}

$w_{K,J}$ est composé de $v_{K,J}$, qui est dans $C$, et $u'_{K,J}$, qui l'est car déduit de $u_{K,J}$ par cochangement de base.

(\textbf{N.B.} On a déjà fait le raisonnement avec un seul $u_\bullet$, et \textit{trois} ouverts $J \subset K \subset L$. Il \msp{43} y aurait lieu d'introduire la notion de $C$-carré
\begin{tikzcd}[baseline=(current bounding box.center)]
X_0 \arrow[r] \arrow[d] & X_1 \arrow[d] \\
X_2 \arrow[r] & X_3
\end{tikzcd}, comme un carré tel que
\printedpage{23}
% head: A. Grothendieck : Dérivateurs, Ch. XVII, Édition des Universités de Montpellier II et Paris VII
$X_1 \amalg_{X_0} X_2 \longrightarrow X_3$ soit dans $C$ (\textbf{N.B.} notion symétrique pour les foncteurs de symétrie des carrés), et dire que le composé de deux $C$-carrés est un $C$-carré.)

Cochangement de base
\[
\begin{tikzcd}
X_\bullet \arrow[r, "u"] \arrow[d] & Y_\bullet \arrow[d] \\
X'_\bullet \arrow[r, "u'"] & Y'_\bullet \rlap{\ .}
\end{tikzcd}
\]

Considérons $K \subset J$, on a
\[
\begin{tikzcd}
& X_K \arrow[r] \arrow[d] \arrow[ddl] \arrow[dr, phantom, "\text{cocart.}"] & X_J \arrow[d] \arrow[dr] & \\
& Y_K \arrow[r] \arrow[ddl] \arrow[dl, phantom, "\text{cocart.}"] & Y_{K,J} \arrow[r, "u_{K,J}"] \arrow[ddl] \arrow[d, phantom, "\text{cocartésien}"] & Y_J \arrow[ddl] \\
X'_K \arrow[r, crossing over] \arrow[d] \arrow[dr, phantom, "\text{cocart.}"] & X'_J \arrow[d] \arrow[dr, crossing over] \arrow[from=1-3, crossing over] & & \\
Y'_K \arrow[r] & Y'_{K,J} \arrow[r, "u'_{K,J}"] & Y'_J \rlap{\ ,} &
\end{tikzcd}
\]
où les trois carrés latéraux marqués “cocart.” sont cocartésiens, ainsi que le carré en plan incliné. Il en résulte, d'abord que le quatrième carré latéral du cube est lui aussi cocartésien, \msp{44} puis que le carré ayant deux côtés $u_{K,J}$, $u'_{K,J}$ est cocartésien. (\textbf{N.B.} comparer le diagramme similaire p. 32 -- il y aurait un lemme commun à dégager, distinguant deux cas \ldots)

$\mathrm{Cof}_2$. Pour prouver que $C_{I^\circ} \cap W_{I^\circ}$ est stable par cochangement de base, je ne vois pas de façon de le faire, si je ne sais déjà qu'on a
\[ C_{I^\circ} \cap W_{I^\circ} \ = \ (TC)_{I^\circ}\ , \tag*{$(*)$} \]
et alors on applique le résultat précédent (où $W$ n'entrait pas) à $TC$ au lieu de $C$. Mais pour avoir $(*)$, il faut que je suppose qu'on est dans l'un des deux cas suivants (de façon à pouvoir appliquer le théorème 2.1 et son corollaire 2.5).
\begin{enumerate}[leftmargin=2.5em]
\item[a)] $I$ fini, ou plus généralement $I$ artinien \textit{et} $\mathrm{Top}(I)$ est noethérien.
\item[b)] $I$ artinien, et de plus $C$ et $TC$ sont stables par composition transfinie.
\end{enumerate}
$\mathrm{Cof}_3$. L'axiome de facorisation me paraît a priori le plus délicat\footnote[15]{L'argument qui suit s'applique à d'autres situations de factorisation, par exemple pour un précouple de Quillen donné $(\Phi, \Psi)$ dans $\M$.}. Mais cela sort finalement sans trop de mal. Soit donné % [sic] facorisation -> factorisation, missing letter
\[ f \ : \ X_\bullet \longrightarrow Y_\bullet \]

\end{document}
